\chapter{结论与可证伪研究议程}
\label{chap:conclusions-research-agenda}

\section{八条结论及其证据位置}

本章不重新复述模型和企业名录，只回答全书核心问题。表\ref{tab:conclusions-evidence-map}中的章节号是结论的主证据位置，不表示其他章节无关。

\begin{table}[H]
\begingroup\hbadness=10000
\centering\small
\caption{结论--证据章节映射}
\label{tab:conclusions-evidence-map}
\begin{tabularx}{\textwidth}{p{0.05\textwidth}Yp{0.2\textwidth}}
\toprule
序 & 可检验结论 & 主证据章节 \\
\midrule
1 & 具身智能是受物理、时序和安全约束的闭环系统，不是语言模型加机械外壳。 & 第 3、6、18--21 章 \\
2 & 接触、状态估计、规划和低层控制仍决定可执行边界；基础模型改变接口与泛化方式，不取消这些约束。 & 第 5--8、14--17 章 \\
3 & 机器人学习的主要瓶颈从单次成功转向数据覆盖、分布偏移、恢复和持续验证。 & 第 9--13、26--28 章 \\
4 & VLA 的差异必须沿动作表示、数据、目标函数、闭环频率和部署接口比较，模型名称本身没有解释力。 & 第 22--25 章 \\
5 & 可交付系统需要感知、状态、规划、技能、控制、安全和运维接口闭合；端到端策略也需承担明确责任边界。 & 第 18--21、29--30 章 \\
6 & 结构化高价值任务会先规模化，因为环境、质量和异常可工程化；开放环境的价值需与人工支持同时核算。 & 第 31--33 章 \\
7 & 企业成熟度由客户级持续运行、制造一致性和运维分母决定，融资、demo、订单和规划产能不可互换。 & 第 34--35 章 \\
8 & 政策和标准改变资源与准入约束，但技术胜负最终仍由可审计性能、单位经济性和责任闭环决定。 & 第 30、36--37 章 \\
\bottomrule
\end{tabularx}
\endgroup
\end{table}

跨本体数据和通用策略已展示迁移可能\cite{openx2024,octo2024}。真实部署材料也开始给出运行与产出分母\cite{figure2025bmw,agility2026totes}；但这些证据尚未共同证明任意环境、长期无人值守和正单位经济性。

\section{保守与激进情景只由触发条件定义}

\textbf{保守情景}被以下事实触发：跨站点成功率不随数据/模型扩大而提高；人工介入与硬件故障使 TCO 长期高于专机或人工；跨本体适配仍需逐任务重采和重验证；法规、保险和责任证据无法随软件更新闭合。此时市场主要停留在结构化制造、物流、危险作业和高价值辅助任务，通用本体更多作为研发平台。

\textbf{激进情景}要求同时出现：多客户系统以版本化口径公开跨季度运行、低介入和安全暴露量；通用预训练显著降低新任务数据与工程工时；量产良率、备件和远程运维使单位成本跨过客户阈值；模型更新可通过自动回归、风险案例和合格评定管理。只满足模型能力或硬件降价之一，不足以触发该情景。

\begin{figure}[H]
\centering
\small
\begin{tabular}{c@{\quad}c@{\quad}c@{\quad}c}
\fbox{跨站点低介入} & \fbox{适配成本下降} & \fbox{制造/运维收敛} & \fbox{更新可验证}
\end{tabular}
\caption{激进情景的联合触发器。四项是“与”关系，不是任选一项。}
\label{fig:aggressive-scenario-triggers}
\end{figure}

\section{可证伪研究议程}

下一阶段研究应优先建设八类公共证据：
\begin{enumerate}[leftmargin=2em]
  \item 以排程小时、动作次数和人机接触小时为分母的统一暴露量；
  \item 含失败、恢复、人工接管和损伤标签的数据集，而非只保留成功轨迹；
  \item 在接触、遮挡和部件退化下的状态估计与安全控制联评；
  \item 控制数据量、模型量和任务覆盖度的缩放实验；
  \item 报告适配数据、重整定和验证工时的跨本体迁移；
  \item 小时到周尺度、允许环境变化的长时程自主基准；
  \item 带功耗、热稳态、最坏时延和网络中断的端侧评测；
  \item 把合格产出、人工分钟、停机、能耗、损伤和 TCO 同时纳入技术评测。
\end{enumerate}

每项都可被否证。例如“持续学习降低维护成本”必须在上线新版本后证明旧任务、安全回归和人工工时不恶化；“跨本体基础模型降低适配成本”必须在多个本体上以同一成功阈值比较数据和工程量。

\section{最终判断}

具身智能已经形成可复现的学习方法、跨机构数据、开源模型和早期客户运行证据，但尚未形成覆盖开放环境、长期自主和统一经济性的通用答案。最有价值的研究不是继续扩大叙事，而是把物理机制、数据血缘、闭环评测、制造运维和责任证据接成同一条可审计链。任何未来判断，只要不能说明何种观测会使其失效，就不应进入本书结论。
